\section{Deep Convection}


Convection transports heat, moisture and momentum vertically.


ACF convection schemes represent:


\begin{itemize}

\item Convective instability

\item Updrafts

\item Downdrafts

\item Latent heat release

\end{itemize}



Important parameters:


\begin{itemize}

\item CAPE

\item CIN

\item Convective precipitation

\item Cloud base height

\end{itemize}



Convective Available Potential Energy (CAPE) is the vertically integrated
positive buoyancy of a lifted parcel between its Level of Free Convection
(LFC) and Equilibrium Level (EL), using virtual temperature $T_v$ to
account for moisture loading (Doswell \& Rasmussen, 1994):

\begin{equation}
\mathrm{CAPE}
=
g\int_{\mathrm{LFC}}^{\mathrm{EL}}
\frac{T_{v,parcel}-T_{v,env}}{T_{v,env}}
\, dz
\end{equation}


Convective Inhibition (CIN) is the corresponding negative-buoyancy integral
below the LFC that must be overcome to initiate convection:

\begin{equation}
\mathrm{CIN}
=
-g\int_{\mathrm{SFC}}^{\mathrm{LFC}}
\frac{T_{v,parcel}-T_{v,env}}{T_{v,env}}
\, dz
\end{equation}


An estimate of the maximum convective updraft speed follows directly from
CAPE (energy conservation for an unentrained parcel):

\begin{equation}
w_{max} = \sqrt{2\,\mathrm{CAPE}}
\end{equation}
